\documentclass[../DoAn.tex]{subfiles}
\begin{document}

\begin{center}
    \Large{\textbf{ABSTRACT}}\\
\end{center}
\vspace{1cm}
Modern IoT and broadband systems combine distributed assets such as CPE devices, home routers, ONTs, ACS/USP management platforms, ISP backends, and SOC/SIEM infrastructures. These assets have long life cycles, limited resources, heterogeneous firmware, and large-scale deployment constraints; therefore, security management must be organized as a continuous process rather than as isolated technical controls. This report studies security management processes for IoT and broadband systems through five pillars: vulnerability assessment, threat modeling, logging and audit, incident response, and patch management. The analysis is based on NIST SP 800-30, NIST SP 800-61, ISO/IEC 27035, MITRE ATT\&CK for IoT, and the research survey files provided in the source folder. The report presents a reference architecture, a data-flow diagram, a vulnerability assessment workflow, an incident response life cycle, and a secure OTA patch management process. The main result is an integrated framework for asset inventory, CVE/CVSS-based prioritization, threat scenario analysis using STRIDE/PASTA/DREAD, security log collection, incident playbook execution, and controlled firmware update deployment. The report also identifies open research gaps, including reliable firmware benchmarks, unified telemetry for embedded devices, legacy devices that cannot be patched, and the need for automation between SOC, NOC, and device management systems.

\end{document}